\lecture{3}{1. September 2025}{Systems of forces and moments}

\begin{problemwithimage}{f3_92}{*3-92}
  Replace the force system by an equivalent resultant force and couple moment at point $O$.
  Take $\Vec{F}_3 = \{ -200 \Vec{i} + 500 \Vec{j} - 300 \Vec{k} \} \,  \unit{\N}$.
\end{problemwithimage}

\begin{derivation}
  The resultant force can be found by summing all of the individual forces as:
  \begin{equation*}
    \Vec{F}_{\mathrm{res}} = \begin{pmatrix}
    \qty{0}{\N}\\
    \qty{0}{\N} \\
    -\qty{300}{\N} \\
    \end{pmatrix} + \begin{pmatrix}
      \qty{0}{\N} \\
    \qty{200}{\N} \\
    \qty{0}{\N} \\
    \end{pmatrix} + \begin{pmatrix}
      - \qty{200}{\N} \\
    \qty{500}{\N} \\
    - \qty{300}{\N} \\
    \end{pmatrix} = \begin{pmatrix}
      - \qty{200}{\N} \\
    \qty{700}{\N} \\
    -\qty{600}{\N} \\
    \end{pmatrix}
  .\end{equation*}
  Each of the three moments can be found as
  \begin{align*}
    \Vec{M}_1 &= \begin{pmatrix}
      \qty{0}{\m} \\
    \qty{2}{\m} \\
    \qty{0}{\m} \\
    \end{pmatrix} \times \begin{pmatrix}
      \qty{0}{\N} \\
    \qty{0}{\N} \\
    - \qty{300}{\N} \\
    \end{pmatrix} = \begin{pmatrix}
      -\qty{600}{\N.\m} \\
    \qty{0}{\N.\m} \\
    \qty{0}{\N.\m} \\
    \end{pmatrix} \\
    \Vec{M}_2 &= \begin{pmatrix}
      \qty{1.5}{\m} \\
    \qty{3.5}{\m} \\
    \qty{0}{\m} \\
    \end{pmatrix} \times \begin{pmatrix}
      \qty{0}{\N} \\
    \qty{200}{\N} \\
    \qty{0}{\N} \\
    \end{pmatrix}  = \begin{pmatrix}
      \qty{0}{\N.\m} \\
    \qty{0}{\N.\m} \\
    \qty{300}{\N.\m} \\
    \end{pmatrix} \\
    \Vec{M}_3 &= \begin{pmatrix}
      \qty{1.5}{\m} \\
    \qty{2}{\m} \\
    \qty{0}{\m} \\
    \end{pmatrix} \times \begin{pmatrix}
      - \qty{200}{\N} \\
    \qty{500}{\N} \\
    - \qty{300}{\N} \\
    \end{pmatrix} = \begin{pmatrix}
      -\qty{600}{\N.\m} \\
    \qty{450}{\N.\m} \\
    \qty{1,15e3}{\N.\m} \\
    \end{pmatrix}
  .\end{align*}
  Summing these three our resultant moment $\Vec{M}_{\mathrm{res}}$ can be found as:
  \begin{equation*}
    \Vec{M}_{\mathrm{res}} = \begin{pmatrix}
      - \qty{600}{\N.\m} \\
    \qty{0}{\N.\m} \\
    \qty{0}{\N.\m} \\
    \end{pmatrix} + \begin{pmatrix}
      \qty{0}{\N.\m} \\
    \qty{0}{\N.\m} \\
    \qty{300}{\N.\m} \\
    \end{pmatrix} + \begin{pmatrix}
      - \qty{600}{\N.\m} \\
    \qty{450}{\N.\m} \\
    \qty{1.15}{\kN.\m} \\
    \end{pmatrix} = \begin{pmatrix}
      - \qty{1200}{\N.\m} \\
    \qty{450}{\N.\m} \\
    \qty{1450}{\N.\m} \\
    \end{pmatrix}
  .\end{equation*}
\end{derivation}

\finalresult{
  \begin{aligned}
    \Vec{F}_{\mathrm{res}} &= \begin{pmatrix}
      - 200 \\
      700 \\
      -600 \\
    \end{pmatrix} \, \unit{\N} \\
    \Vec{M}_{\mathrm{res}} &= \begin{pmatrix}
      -1200 \\
      450\\
      1450\\
    \end{pmatrix} \, \unit{\N.\m}
  .\end{aligned}
}

\begin{problemwithimage}{f3_134}{3-134}
  Replace the loading by a single resultant force, and specify the location of the force measured from point $O$.
\end{problemwithimage}

\begin{derivation}
  First we will find the position of the resultant of the distributed loading.
  To do this, we split the distributed loading into two parts $P_1$ and $P_2$ split at the top point of the distributed loading.
  We get
  \begin{align*}
    \left( P_1 \right)_{\mathrm{res}} &= \frac{1}{2}\cdot \qty{7.5}{\m} \cdot \left( - \qty{6}{\kN\per\m} \right) = - \qty{225}{\kN} \\
    \left( P_2 \right)_{\mathrm{res}} &= \frac{1}{2} \cdot \qty{4,5}{\m} \cdot \left( - \qty{6}{\kN\per\m}  \right) = -\qty{135}{\kN} 
  .\end{align*}
  The total resultant force is therefore:
  \begin{equation*}
    F_{\mathrm{res}} = -\qty{225}{\kN} - \qty{135}{\kN} - \qty{15}{\kN}  = \qty{-375}{\kN}
  .\end{equation*}
  
  To find the location of this force we must find the position to place this force at to give the same couple moment. 
  We add all the individual moments:
  \begin{equation*}
    M_{\mathrm{res}} = -\qty{500}{\kN.\m} + \left( \qty{5}{\m} \cdot (- \qty{225}{\kN})  \right) + \left( \qty{9}{\m} \cdot (-\qty{135}{\kN})  \right) + \left( \qty{12}{\m} \cdot(- \qty{15}{\kN})  \right) = -\qty{3020}{\kN.\m} 
  .\end{equation*}
  Now we divide this by the magnitude of the force to find the required distance from $O$:
  \begin{equation*}
    x_{F_{\mathrm{res}}} = \frac{\qty{3020}{\kN.\m} }{\qty{375}{\kN} } = \qty{8.05}{\meter}
  .\end{equation*}
\end{derivation}

\finalresult{
\begin{aligned}
  F_{\mathrm{res}} &= - \qty{375}{\kN} \\
  x_{F_{\mathrm{res}}} &= \qty{8,05}{\m} 
.\end{aligned}
}
